% Generated draft: requires mathematical and source review.
\documentclass[11pt]{article}
\usepackage[margin=25mm]{geometry}
\usepackage{fontspec}
\setmainfont{Latin Modern Roman}
\usepackage{amsmath,amssymb,mathtools,unicode-math}
\setmathfont{Latin Modern Math}
\usepackage{enumitem,needspace,xurl,hyperref}
\setlength{\parindent}{0pt}
\setlength{\parskip}{5pt}
\setlength{\emergencystretch}{3em}
\newcommand{\N}{\mathbb{N}}
\newcommand{\Z}{\mathbb{Z}}
\newcommand{\Q}{\mathbb{Q}}
\newcommand{\R}{\mathbb{R}}
\newcommand{\C}{\mathbb{C}}
\newcommand{\F}{\mathbb{F}}
\newcommand{\A}{\mathbb{A}}
\newcommand{\PP}{\mathbb{P}}
\newcommand{\E}{\mathbb{E}}
\newcommand{\eps}{\varepsilon}
\newcommand{\dd}{\,\mathrm{d}}
\begin{document}
\section*{UnsolvedMath ID: 30002828}
\textbf{Source problem code:} catalogue entry
\subsection{Definitions and mathematical statement}
Let \(Y\) be a maximally degenerating projective conifold with mirror \(\check Y\), as in the supplied source, and let \(S\subset Y\) and \(\check S\subset\check Y\) be their respective sets of conifold points. Put \(U=Y\setminus S\) and \(\check U=\check Y\setminus\check S\). The source does not specify the dimension, the precise meaning of maximal degeneration, or the construction and hypotheses of the mirror correspondence; these remain qualifications on the statement below.

Write \(\Omega_Y^2\) and \(\Omega_{\check Y}^1\) for the sheaves intended by the source's notation \(\Omega^2\) and \(\Omega^1\), respectively. Their definitions at the singular points are not supplied. Once these sheaves and their restrictions to the complements are fixed, the localization sequences give connecting maps
\[
\delta_Y:H^1(U,\Omega_Y^2|_U)\longrightarrow H^2_S(Y,\Omega_Y^2),\qquad
\delta_{\check Y}:H^1(\check U,\Omega_{\check Y}^1|_{\check U})\longrightarrow H^2_{\check S}(\check Y,\Omega_{\check Y}^1).
\]
Here \(H^i_T(X,\mathcal F)\) denotes sheaf cohomology with support in \(T\). The source supplies identifications
\[
H^2_S(Y,\Omega_Y^2)\cong\C^S,\qquad
H^2_{\check S}(\check Y,\Omega_{\check Y}^1)\cong\C^{\check S}.
\]
Use these intended identifications to denote the coordinates of the connecting-map images by \((\delta_Y\alpha)_s\) and \((\delta_{\check Y}\beta)_{\check s}\). The source does not specify the sheaf conventions or hypotheses establishing these identifications.

The proposed equivalence asks whether, for every pair satisfying the intended source hypotheses, the following four conditions are equivalent:
\begin{enumerate}
\item \(Y\) admits an analytic smoothing, understood as a flat holomorphic deformation over a pointed complex disc with central fiber isomorphic to \(Y\) and nonsingular fibers over sufficiently small nonzero parameters. Additional requirements on the family, such as properness, are not specified in the record.
\item There exists \(\alpha\in H^1(U,\Omega_Y^2|_U)\) such that
\[
\forall s\in S,\qquad (\delta_Y\alpha)_s\ne0.
\]
\item There exists \(\beta\in H^1(\check U,\Omega_{\check Y}^1|_{\check U})\) such that
\[
\forall\check s\in\check S,\qquad (\delta_{\check Y}\beta)_{\check s}\ne0.
\]
\item \(\check Y\) admits a projective resolution. The source does not specify whether projectivity refers to the resolution morphism or its total space, or whether a small or crepant resolution is required.
\end{enumerate}
The missing conventions above must be resolved before this becomes a fully specified mathematical assertion. The supplied dated assessment says that the equivalences and a mirror conifold transition are proved in the formal Gross--Siebert setting, while lifting the families to analytic families was explicitly still ongoing in the source. This is supplied curation, not an independently verified theorem or current status assessment.
\subsection{Short English statement}
For the intended class of maximally degenerating projective conifolds and their mirrors, ask whether analytic smoothability of \(Y\), the existence of cohomology classes with nonzero local coordinates at every conifold point on each side, and existence of a projective resolution of \(\check Y\) are equivalent. Precise degeneration, mirror, differential-sheaf, and resolution conventions remain unspecified. The supplied partially solved label concerns formal results and an unresolved analytic lifting qualification.
\subsection{Sources}
\begin{itemize}
\item Local UnsolvedMath registry: supplied complete source record, permanent ID \(30002828\), title \emph{Mirror Conifold Transitions and Analytic Smoothings}, exact source problem code \texttt{catalogue entry}, in the Oberwolfach Reports Open Problems collection. The project identifies \texttt{data/unsolvedmath/problems.json} and \texttt{catalog/problems.json} as registry files; the supplied record does not identify which contains this entry.
\item Supplied report citation: \emph{Algebraic Geometry} (2015), Oberwolfach Reports, pp.\ 783--836, DOI \texttt{10.4171/owr/2015/15}. \url{https://doi.org/10.4171/owr/2015/15}. The record describes this collection as extracted from workshop open-problem discussions and supplies proposed year 2015, but names no proposer and gives no problem-specific page within the report.
\item Supplied literature source: Gross--Siebert, \emph{affine geometry and mirror symmetry}. \url{https://arxiv.org/abs/math/0309070}. The record supplies no theorem or section locator connecting this reference to the precise equivalence.
\item Literature assessment dated exactly \texttt{2026-08-21}: supplied status \texttt{partially\_solved}. The supplied assessment attributes the equivalences and mirror conifold transition to the formal Gross--Siebert setting, and says lifting the families to analytic families was explicitly still ongoing in the source. This is supplied dated curation derived from the OWR record, not independently verified here; its time-sensitive status and cited claims require independent verification.
\item The registry labels the statement \texttt{corrected\_verified} and reports manual semantic, typographic, and LaTeX review of a meaning-preserving cleanup. These are supplied editorial claims, not independent mathematical verification of this draft.
\end{itemize}
\end{document}
